% Deutsche Version
\documentclass[12pt,a4paper]{article}
\usepackage[utf8]{inputenc}
\usepackage[T1]{fontenc}
\usepackage{lmodern}
\usepackage{amsmath,amssymb}
\usepackage{graphicx}
\usepackage{xcolor}
\usepackage{hyperref}
\usepackage{geometry}
\geometry{a4paper, left=3cm, right=3cm, top=3cm, bottom=3cm}
\usepackage{setspace}
\onehalfspacing
\usepackage{parskip}
\usepackage[ngerman]{babel}
\usepackage{csquotes}
\usepackage{microtype}
\usepackage{booktabs}
\usepackage{longtable}
\usepackage{array}
\usepackage{listings}
\usepackage{caption}
\usepackage{subcaption}
\usepackage{float}
\usepackage{url}
\usepackage{natbib}
\usepackage{titling}

\lstset{
  language=Python,
  basicstyle=\ttfamily\small,
  keywordstyle=\color{blue},
  commentstyle=\color{green!40!black},
  stringstyle=\color{red},
  showstringspaces=false,
  numbers=left,
  numberstyle=\tiny,
  numbersep=5pt,
  breaklines=true,
  frame=single,
  backgroundcolor=\color{gray!5},
  tabsize=2,
  captionpos=b
}

\title{\Huge\textbf{Von Scheme zu DeepProbLog} \\[2mm]
       \LARGE Die ARS als methodologischer Bauplan \\[2mm]
       \LARGE für moderne neuro-symbolische Programmierung}

\author{
  \large
  \begin{tabular}{c}
    Paul Koop
  \end{tabular}
}

\date{\large 1994--2026}

\begin{document}

\maketitle

\begin{abstract}
Dieser Beitrag zeichnet die methodologische Kontinuität von frühen Implementierungen 
der Algorithmisch Rekursiven Sequenzanalyse (ARS) in Scheme, Pascal und Lisp 
(1992--1994) zu zeitgenössischen neuro-symbolischen Programmierframeworks wie 
DeepProbLog (2018) nach. Ich argumentiere, dass die ARS bereits die Kernprinzipien 
neuro-symbolischer Integration – Mustererkennung (System 1), regelbasiertes 
Schließen (System 2), probabilistische Unsicherheitsquantifizierung und Erklärbarkeit 
durch Design – Jahrzehnte vor der Prägung des Begriffs "neuro-symbolische KI" 
verkörperte. Der Beitrag rekonstruiert zunächst die proto-neuro-symbolische 
Architektur der ARS, stellt dann DeepProbLog als modernes Framework vor, das 
ähnliche Prinzipien mit größerer Skalierbarkeit implementiert, und demonstriert 
schließlich eine DeepProbLog-Implementierung des klassischen ARS-Verkaufsgesprächs-Korpus. 
Die Synthese zeigt, dass die ARS einen methodologischen Bauplan liefert, den 
DeepProbLog technisch instanziieren kann. Ich schließe mit einer Forschungsagenda 
für die Integration der methodologischen Strenge der ARS mit der Rechenleistung 
von DeepProbLog.
\end{abstract}

\newpage
\tableofcontents
\newpage

\section{Einleitung: Von Lisp zu DeepProbLog}

Die Algorithmisch Rekursive Sequenzanalyse (ARS), dokumentiert in den frühen 
Jupyter-Notebooks und Codedateien von 1992--1994, stellt einen der frühesten 
systematischen Versuche dar, Mustererkennung mit regelbasiertem Schließen in 
der Analyse sequenzieller sozialer Interaktionen zu integrieren. Die drei 
Kernimplementierungen –

\begin{itemize}
    \item \textbf{Induktor in Scheme}: Induktion probabilistischer kontextfreier 
    Grammatiken (PCFG) aus Terminalzeichenketten durch Übergangszählung,
    \item \textbf{Parser in Pascal}: Validierung der Wohlgeformtheit von 
    Sequenzen mittels eines Chart-Parsers,
    \item \textbf{Transduktor in Lisp}: Generierung neuer Sequenzen aus der 
    induzierten Grammatik,
\end{itemize}

– verkörpern gemeinsam das, was heute als \textbf{neuro-symbolische KI} bezeichnet 
wird. Sie kombinieren datengetriebene Musterentdeckung (der Induktor) mit 
symbolischer Regelanwendung (der Parser und Transduktor) und quantifizieren 
Unsicherheit durch probabilistische Gewichte.

In den drei dazwischenliegenden Jahrzehnten hat das Feld anspruchsvollere 
Frameworks für neuro-symbolische Integration entwickelt. Eines der prominentesten 
ist \textbf{DeepProbLog} \citep{manhaeve2018deepproblog}, das die probabilistische 
Logikprogrammiersprache ProbLog um neuronale Prädikate erweitert, die von 
tiefen Netzwerken gelernt werden. DeepProbLog erlaubt es Benutzern, symbolische 
Regeln mit Wahrscheinlichkeiten zu definieren, während neuronale Netze die 
Wahrscheinlichkeiten von Grundtatsachen aus Daten lernen.

Dieser Beitrag leistet drei Dinge:

\begin{enumerate}
    \item Er rekonstruiert die ARS-Architektur als \textbf{proto-neuro-symbolisches} 
    System und bildet ihre Komponenten auf zeitgenössische neuro-symbolische 
    Konzepte ab.
    
    \item Er stellt DeepProbLog als modernes Framework vor, das dieselben 
    Prinzipien mit größerer Skalierbarkeit und neuronaler Integration 
    implementiert.
    
    \item Er präsentiert eine \textbf{DeepProbLog-Implementierung} des 
    klassischen ARS-Verkaufsgesprächs-Korpus und demonstriert, wie die 
    ARS-Methodologie in ein modernes neuro-symbolisches Framework portiert 
    werden kann.
\end{enumerate}

Die übergreifende These ist, dass \textbf{die ARS einen methodologischen Bauplan 
liefert, den DeepProbLog technisch instanziieren kann}. Die beiden Ansätze sind 
keine Konkurrenten, sondern Komplemente: Die ARS trägt methodologische Strenge 
und interpretative Verankerung bei; DeepProbLog trägt Skalierbarkeit und 
neuronales Lernen bei.

\section{Die ARS-Architektur als proto-neuro-symbolisches System}

\subsection{Drei Komponenten, drei kognitive Funktionen}

Die drei ARS-Implementierungen lassen sich auf die von Kahneman 
\citep{kahneman2011thinking} popularisierte und von der neuro-symbolischen 
KI-Forschung übernommene Unterscheidung zwischen System 1 und System 2 
\citep{marcus2020next} abbilden:

\begin{table}[H]
\centering
\caption{ARS-Komponenten als kognitive Systeme}
\label{tab:cognitive}
\begin{tabular}{@{} p{3cm} p{4cm} p{6cm} @{}}
\toprule
\textbf{Komponente} & \textbf{Kognitive Funktion} & \textbf{Neuro-symbolische Abbildung} \\
\midrule
Induktor (Scheme) & Mustererkennung, Übergangszählung & System 1 (Lernen aus Daten) \\
Parser (Pascal) & Strukturelle Validierung, Wohlgeformtheitsprüfung & System 2 (Regelanwendung) \\
Transduktor (Lisp) & Generative Regelanwendung & System 2 (symbolische Generierung) \\
Multiagent (Python) & Rollenzuweisung, Interaktion & Hybrid (Entscheidungsbaum + Grammatik) \\
\bottomrule
\end{tabular}
\end{table}

\subsection{Die probabilistische Grammatik als neuro-symbolische Schnittstelle}

Die induzierte probabilistische kontextfreie Grammatik (PCFG) dient als zentrale 
neuro-symbolische Schnittstelle:

\begin{verbatim}
(KBG -> . VBG)
(VBG -> . KBBd)
(KBBd -> . VBBd)
(VBBd -> . KBA)
(KBA -> . VBA)
(VBA -> . KBBd) (VBA -> . KAE)
(KAE -> . VAE)
(VAE -> . KAE) (VAE -> . KAA)
(KAA -> . VAA)
(VAA -> . KAV)
(KAV -> . VAV)
\end{verbatim}

Jede Produktionsregel hat eine Wahrscheinlichkeit (in dieser vereinfachten 
Grammatik implizit 1.0, in der vollständigen Implementierung gewichtet nach 
empirischen Häufigkeiten). Die Grammatik ist gleichzeitig:

\begin{itemize}
    \item \textbf{Symbolisch}: Regeln sind explizit, inspizierbar und 
    falsifizierbar.
    \item \textbf{Probabilistisch}: Regelanwendungen haben Wahrscheinlichkeiten 
    basierend auf empirischen Häufigkeiten.
    \item \textbf{Generativ}: Neue Sequenzen können durch Anwendung der Regeln 
    erzeugt werden.
    \item \textbf{Verifizierbar}: Der Parser kann prüfen, ob eine Sequenz 
    wohlgeformt ist.
\end{itemize}

Diese vier Eigenschaften sind genau das, was zeitgenössische neuro-symbolische 
Frameworks anstreben.

\subsection{Das Multiagentensystem als neuro-symbolischer Prototyp}

Das Python-Multiagentensystem (Zellen 29--33 im Notebook) ist besonders 
aufschlussreich:

\begin{lstlisting}[caption=Multiagenten-Rollenzuweisung]
# Entscheidung über die Rollenverteilung basierend auf Ware und Zahlungsmittel
if agent_k_ware > agent_v_ware:
    agent_k_role = 'Käufer'
    agent_v_role = 'Verkäufer'
else:
    agent_k_role = 'Verkäufer'
    agent_v_role = 'Käufer'
\end{lstlisting}

Dieser Entscheidungsbaum ist eine \textbf{symbolische Regel} (System 2), die 
Agentenrollen basierend auf einem einfachen Muster (System 1: Vergleich zweier 
Zahlen) bestimmt. Die anschließende Interaktion folgt der probabilistischen 
Grammatik. Dies ist eine hybride Architektur: Die Rollenzuweisung ist 
deterministisch und regelbasiert; die Dialoggenerierung ist probabilistisch 
und grammatikbasiert.

Die ARS antizipiert damit das \textbf{Neural | Symbolic}-Muster in Kautz' 
Taxonomie \citep{kautz2020third}: neuronale (oder heuristische) Wahrnehmung 
bestimmt symbolische Rollen; symbolisches Schließen (die Grammatik) steuert 
das anschließende Verhalten.

\section{DeepProbLog: Ein modernes neuro-symbolisches Framework}

\subsection{Was DeepProbLog ist}

DeepProbLog \citep{manhaeve2018deepproblog} erweitert die probabilistische 
Logikprogrammiersprache ProbLog um \textbf{neuronale Prädikate}. Ein neuronales 
Prädikat ist ein Prädikat, dessen Wahrheitswahrscheinlichkeit von einem 
neuronalen Netz berechnet wird. Ein neuronales Prädikat `digit(image, d)` könnte 
beispielsweise die Wahrscheinlichkeit repräsentieren, dass ein Bild die Ziffer 
`d` zeigt.

DeepProbLog-Programme bestehen aus:

\begin{itemize}
    \item \textbf{Fakten}: Grundatome mit Wahrscheinlichkeiten (z.B. `0.5::edge(a,b)`).
    \item \textbf{Regeln}: Logische Implikationen (z.B. `path(X,Y) :- edge(X,Y)`).
    \item \textbf{Neuronale Prädikate}: Durch neuronale Netze definierte Prädikate.
    \item \textbf{Anfragen}: Zu beantwortende probabilistische Fragen.
\end{itemize}

Die Inferenz in DeepProbLog berechnet die Wahrscheinlichkeit einer Anfrage 
gegeben das Programm und die Ausgaben des neuronalen Netzes. Lernen aktualisiert 
die neuronalen Netzgewichte, um die Likelihood der beobachteten Daten zu 
maximieren.

\subsection{Abbildung von ARS-Konzepten auf DeepProbLog}

\begin{table}[H]
\centering
\caption{Abbildung von ARS auf DeepProbLog}
\label{tab:mapping}
\begin{tabular}{@{} p{4cm} p{4cm} p{4cm} @{}}
\toprule
\textbf{ARS-Konzept} & \textbf{DeepProbLog-Konzept} & \textbf{Erklärung} \\
\midrule
Terminalzeichen & Grundfakten & `KBG`, `VBG`, `KBBd`, etc. \\
Produktionsregeln & Logische Regeln & `next(X,Y) :- transition(X,Y)` \\
Übergangswahrscheinlichkeiten & Faktenwahrscheinlichkeiten & `0.8::next(KBG, VBG)` \\
Induktor (Übergangszählung) & Neuronales Prädikatenlernen & Aus Daten gelernt \\
Parser (Wohlgeformtheit) & Beweissuche & Anfrage `next(KBG, VBG)` \\
Transduktor (Generierung) & Sampling aus Verteilung & `sample(next(Start, X))` \\
Multiagentenrollen & Probabilistische Entscheidungsregeln & Rollenzuweisung mit Wahrscheinlichkeit \\
\bottomrule
\end{tabular}
\end{table}

\subsection{Warum DeepProbLog ein natürlicher Nachfolger der ARS ist}

DeepProbLog bewahrt die wichtigsten methodologischen Tugenden der ARS:

\begin{enumerate}
    \item \textbf{Erklärbarkeit}: Regeln sind explizit und inspizierbar.
    \item \textbf{Probabilistische Unsicherheit}: Wahrscheinlichkeiten 
    quantifizieren Unsicherheit.
    \item \textbf{Generative Kapazität}: Neue Sequenzen können erzeugt werden.
    \item \textbf{Verifizierbarkeit}: Anfragen können geprüft werden.
\end{enumerate}

Aber es fügt Fähigkeiten hinzu, die der ARS fehlen:

\begin{enumerate}
    \item \textbf{Neuronale Integration}: Neuronale Netze können Wahrscheinlichkeiten 
    aus Rohdaten (Bildern, Text, Audio) lernen, nicht nur aus vorkodierten 
    Kategorien.
    
    \item \textbf{Skalierbarkeit}: DeepProbLog kann große Datensätze durch 
    stochastischen Gradientenabstieg verarbeiten.
    
    \item \textbf{Kontinuierliches Lernen}: Das neuronale Netz kann inkrementell 
    aktualisiert werden, wenn neue Daten eintreffen.
    
    \item \textbf{Tiefes Merkmalslernen}: Neuronale Netze können automatisch 
    relevante Merkmale entdecken und reduzieren so den Bedarf an manueller 
    Kategorienbildung.
\end{enumerate}

\section{DeepProbLog-Implementierung des ARS-Korpus}

\subsection{Die Terminalzeichen als probabilistische Fakten}

Der erste Schritt ist die Kodierung der ARS-Terminalzeichen als probabilistische 
Fakten. Die Übergangswahrscheinlichkeiten werden aus dem Korpus gelernt:

\begin{lstlisting}[caption=DeepProbLog-Kodierung der ARS-Grammatik]
% DeepProbLog-Implementierung der ARS-Grammatik für Verkaufsgespräche
% Basierend auf dem Aachener Markt-Transkript (1994)

% Terminalzeichen als Prädikate
predicate(kbg/0). predicate(vbg/0). predicate(kbbd/0). predicate(vbbd/0).
predicate(kba/0). predicate(vba/0). predicate(kae/0). predicate(vae/0).
predicate(kaa/0). predicate(vaa/0). predicate(kav/0). predicate(vav/0).

% Neuronale Prädikate für Übergangswahrscheinlichkeiten
nn(transition, [in:symbol, out:symbol]) :: neural_network.

% Regeln: Wohlgeformte Sequenzen folgen Übergängen
% Startsymbol ist KBG (Kundenbegrüßung)
next(S) :- transition(start, S).

% Rekursive Regel für Sequenzen der Länge > 1
next([A,B|Rest]) :-
    transition(A, B),
    next([B|Rest]).

% Anfrage: Wahrscheinlichkeit, dass eine gegebene Sequenz wohlgeformt ist
query(well_formed(Sequence)) :- next(Sequence).

% Generierung: Sample einer wohlgeformten Sequenz
sample(well_formed(S)) :- next(S).
\end{lstlisting}

\subsection{Lernen der Übergangswahrscheinlichkeiten aus Daten}

Das neuronale Netz für Übergangswahrscheinlichkeiten kann auf den aus dem 
ARS-Korpus extrahierten Terminalzeichensequenzen trainiert werden. Das Korpus 
lautet:

\begin{verbatim}
KBG VBG KBBd VBBd KBA VBA KBBd VBBd KBA VBA KAE VAE KAE VAE KAA VAA KAV VAV
\end{verbatim}

In DeepProbLog können wir dies als Trainingsdaten kodieren:

\begin{lstlisting}[caption=Kodierung der Trainingsdaten]
% Trainingsbeispiele: beobachtete Übergänge
train(transition(kbg, vbg), true).
train(transition(vbg, kbbd), true).
train(transition(kbbd, vbbd), true).
train(transition(vbbd, kba), true).
train(transition(kba, vba), true).
train(transition(vba, kbbd), true).
train(transition(vba, kae), true).
train(transition(kae, vae), true).
train(transition(vae, kae), true).
train(transition(vae, kaa), true).
train(transition(kaa, vaa), true).
train(transition(vaa, kav), true).
train(transition(kav, vav), true).

% Negative Beispiele (optional)
train(transition(kbg, kbbd), false).
train(transition(vbg, vbg), false).
\end{lstlisting}

Das neuronale Netz lernt, beobachteten Übergängen hohe Wahrscheinlichkeiten 
und unbeobachteten niedrige Wahrscheinlichkeiten zuzuweisen. Nach dem Training 
approximiert das Netz die empirischen Übergangshäufigkeiten.

\subsection{Das Multiagentensystem in DeepProbLog}

Das Multiagentensystem kann als eine Menge probabilistischer Regeln mit 
Rollenzuweisung implementiert werden:

\begin{lstlisting}[caption=Multiagentensystem in DeepProbLog]
% Agentenrollen basierend auf Waren- und Geldausstattung
% Diese könnten durch neuronale Netze aus Daten gelernt werden
nn(role, [in:goods, in:money, out:role]) :: role_network.

% Rollenzuweisungsregel
agent_role(A, buyer) :- goods(A, G), money(A, M), role(G, M, buyer).
agent_role(A, seller) :- goods(A, G), money(A, M), role(G, M, seller).

% Interaktionsregeln basierend auf Rollen
utterance(A, kb) :- agent_role(A, buyer), start_turn.
utterance(A, vg) :- agent_role(A, seller), previous_utterance(_, kb).

% Grammatikbasierte Dialogfortsetzung
next_utterance(A, Sym) :-
    previous_utterance(_, PrevSym),
    transition(PrevSym, Sym),
    agent_role(A, _).

% Anfrage: Wahrscheinlichkeitsverteilung der nächsten Äußerung
query(next_utterance(seller, Sym)).
\end{lstlisting}

\subsection{Erklärbarkeit in DeepProbLog}

Ein entscheidender Vorteil von DeepProbLog gegenüber rein neuronalen Netzen 
ist die \textbf{Erklärbarkeit durch Design}. Für jede Anfrage kann DeepProbLog 
einen Beweisbaum liefern:

\begin{lstlisting}[caption=Erklärbarkeitsausgabe]
| ?- explain(well_formed([KBG, VBG, KBBd, VBBd, KBA])).

Beweis:
1. well_formed([KBG, VBG, KBBd, VBBd, KBA])
   ← next([KBG, VBG, KBBd, VBBd, KBA])
2. next([KBG, VBG, KBBd, VBBd, KBA])
   ← transition(KBG, VBG) ∧ next([VBG, KBBd, VBBd, KBA])
3. transition(KBG, VBG) ← neural_network(KBG, VBG) [p = 0.67]
4. next([VBG, KBBd, VBBd, KBA])
   ← transition(VBG, KBBd) ∧ next([KBBd, VBBd, KBA])
5. transition(VBG, KBBd) ← neural_network(VBG, KBBd) [p = 1.00]
   ... (Fortsetzung)

Wahrscheinlichkeit: 0.67 × 1.00 × 0.67 × ... = 0.42
\end{lstlisting}

Dieser Beweisbaum ist direkt interpretierbar und bildet exakt die ARS-Grammatikregeln 
ab. Der einzige Unterschied ist, dass die Wahrscheinlichkeiten von einem 
neuronalen Netz gelernt und nicht manuell gezählt werden.

\subsection{Vergleich mit der ursprünglichen ARS-Implementierung}

\begin{table}[H]
\centering
\caption{ARS vs. DeepProbLog-Implementierung}
\label{tab:comparison}
\begin{tabular}{@{} p{4cm} p{4cm} p{4cm} @{}}
\toprule
\textbf{Kriterium} & \textbf{ARS (Scheme/Lisp)} & \textbf{DeepProbLog} \\
\midrule
Wahrscheinlichkeitslernen & Manuelles Zählen & Neuronales Netzlernen \\
Regelrepräsentation & Assoziationslisten & Logische Prädikate \\
Parsing-Algorithmus & Chart-Parser (handkodiert) & Beweissuche (eingebaut) \\
Generierung & Benutzerdefinierter Transduktor & Sampling aus Verteilung \\
Erklärbarkeit & Nachvollziehbar via Code & Beweisbäume \\
Skalierbarkeit & Gering (n=8) & Hoch (n > 1000) \\
Neuronale Integration & Keine & Vollständig (neuronale Prädikate) \\
\bottomrule
\end{tabular}
\end{table}

Die DeepProbLog-Implementierung bewahrt die methodologischen Tugenden der ARS 
und fügt Skalierbarkeit und neuronales Lernen hinzu. Sie ist kein Ersatz, 
sondern eine \textbf{technische Instanziierung} derselben methodologischen 
Prinzipien.

\section{Auf dem Weg zu einer Synthese: ARS als Bauplan, DeepProbLog als Motor}

\subsection{Was die ARS zu DeepProbLog beiträgt}

Die ARS-Methodologie bietet drei Lehren für DeepProbLog-Praktiker:

\begin{enumerate}
    \item \textbf{Interpretative Verankerung}: Die Bedeutung von Symbolen muss 
    dokumentiert sein. Ein DeepProbLog-Programm mit nicht interpretierten 
    Symbolen ist nicht erklärend. Die ARS zeigt, wie Symbole in qualitativer 
    Interpretation verankert werden können.
    
    \item \textbf{Trennung von Struktur und Statistik}: Die ARS hält eine 
    strikte Trennung zwischen strukturellen Regeln (deterministisch, logisch) 
    und statistischen Regularitäten (probabilistisch, empirisch) ein. DeepProbLogs 
    Mischung aus logischen Regeln und neuronalen Wahrscheinlichkeiten riskiert 
    die Vermischung dieser Ebenen. Die ARS legt nahe, sie in der Programmstruktur 
    getrennt zu halten.
    
    \item \textbf{Falsifizierbarkeit als Validierung}: Die ARS besteht darauf, 
    dass Grammatiken durch Gegenbeispiele falsifizierbar sein müssen. Die 
    Validierung in DeepProbLog stützt sich typischerweise auf Likelihood-Maximierung. 
    Die ARS schlägt vor, dies durch qualitative Falsifikationstests zu ergänzen.
\end{enumerate}

\subsection{Was DeepProbLog zur ARS beiträgt}

Umgekehrt bietet DeepProbLog drei Verbesserungen für ARS-Praktiker:

\begin{enumerate}
    \item \textbf{Skalierbares Lernen}: Das manuelle Übergangszählen der ARS 
    skaliert nicht. DeepProbLogs neuronales Lernen kann tausende Beispiele 
    verarbeiten.
    
    \item \textbf{Rohdatenintegration}: Die ARS benötigt vorkodierte Terminalzeichen. 
    DeepProbLog kann direkt aus Rohdaten (Text, Bilder, Audio) durch neuronale 
    Prädikate lernen.
    
    \item \textbf{Kontinuierliche Aktualisierung}: ARS-Grammatiken sind statisch. 
    DeepProbLog-Netze können inkrementell aktualisiert werden, wenn neue Daten 
    eintreffen.
\end{enumerate}

\subsection{Eine Forschungsagenda für neuro-symbolische ARS}

Basierend auf dieser Synthese schlage ich eine Forschungsagenda vor:

\begin{enumerate}
    \item \textbf{Portierung des ARS-Korpus nach DeepProbLog}: Vervollständigung 
    der Implementierung der Verkaufsgesprächs-Grammatik in DeepProbLog, 
    einschließlich aller 12 Terminalzeichen und ihrer Übergangswahrscheinlichkeiten.
    
    \item \textbf{Hinzufügung neuronaler Prädikate für Rohdaten}: Training 
    neuronaler Netze zur direkten Abbildung von Rohdaten (Audio, Transkript) 
    auf Terminalzeichen.
    
    \item \textbf{Implementierung des Multiagentensystems}: Aufbau eines 
    vollständigen Multiagentensystems, in dem Agenten Rollen und Interaktionsmuster 
    durch DeepProbLog lernen.
    
    \item \textbf{Validierung mit den XAI-Kriterien}: Evaluation der DeepProbLog-Implementierung 
    anhand der XAI-Kriterien der ARS (Verständlichkeit, Genauigkeit, Wissensgrenzen).
    
    \item \textbf{Skalierung auf größere Korpora}: Anwendung der DeepProbLog-ARS 
    auf größere Datensätze (hunderte oder tausende Gespräche) zur Testung der 
    Skalierbarkeit.
\end{enumerate}

\section{Fazit}

Dieser Beitrag hat die methodologische Kontinuität von frühen ARS-Implementierungen 
in Scheme, Pascal und Lisp zu zeitgenössischer neuro-symbolischer Programmierung 
in DeepProbLog nachgezeichnet. Ich habe argumentiert, dass die ARS bereits die 
Kernprinzipien neuro-symbolischer Integration – Mustererkennung, regelbasiertes 
Schließen, probabilistische Unsicherheit und Erklärbarkeit durch Design – 
Jahrzehnte vor der Prägung des Begriffs verkörperte.

Die Abbildung von ARS-Konzepten auf DeepProbLog ist direkt und natürlich. Die 
probabilistische Grammatik wird zu einer Menge logischer Regeln mit neuronalen 
Prädikaten; der Parser wird zur Beweissuche; der Transduktor wird zum Sampling. 
DeepProbLog ersetzt die ARS nicht, sondern \textit{instanziiert} ihren 
methodologischen Bauplan mit modernen Rechenwerkzeugen.

Die Synthese ist kein Wettbewerb, sondern eine Komplementarität. Die ARS liefert 
die methodologische Strenge und interpretative Verankerung, die DeepProbLog 
(und der neuro-symbolischen KI allgemein) oft fehlt. DeepProbLog liefert die 
Skalierbarkeit und das neuronale Lernen, das der ARS fehlt. Zusammen deuten 
sie auf ein \textbf{methodologisch fundiertes, skalierbares neuro-symbolisches 
Framework} für die Analyse sequenzieller sozialer Interaktionen hin.

Die Frage für zukünftige Forschung ist nicht, ob die ARS oder DeepProbLog 
überlegen ist. Beide sind Werkzeuge für unterschiedliche Zwecke. Die Frage 
ist, wie sie integriert werden können, so dass die methodologischen Lehren 
der ARS die technische Entwicklung von DeepProbLog informieren und die 
Rechenleistung von DeepProbLog die Reichweite der ARS erweitert.

\newpage
\begin{thebibliography}{99}

\bibitem[Kahneman(2011)]{kahneman2011thinking}
Kahneman, D. (2011). \textit{Thinking, Fast and Slow}. Farrar, Straus and Giroux.

\bibitem[Kautz(2020)]{kautz2020third}
Kautz, H. (2020). The third AI summer: AAAI Robert S. Engelmore Memorial Award 
Lecture. \textit{AI Magazine}, 43(1), 93-104.

\bibitem[Koop(1992)]{koop1992parser}
Koop, P. (1992). \textit{Demo-Parser Chart-Parser Version 1.0}. Pascal-Quellcode.

\bibitem[Koop(1994)]{koop1994scheme}
Koop, P. (1994). \textit{Grammatikinduktion empirisch gesicherter 
Verkaufsgespräche}. Scheme-Quellcode.

\bibitem[Koop(1994)]{koop1994lisp}
Koop, P. (1994). \textit{Sequenzanalyse empirisch gesicherter 
Verkaufsgespräche}. Lisp-Quellcode.

\bibitem[Koop(2023)]{koop2023notebook}
Koop, P. (2023). \textit{Qualitative Sozialforschung und Große Sprachmodelle}. 
Jupyter Notebook.

\bibitem[Manhaeve et al.(2018)]{manhaeve2018deepproblog}
Manhaeve, R., Dumancic, S., Kimmig, A., Demeester, T., \& De Raedt, L. (2018). 
DeepProbLog: Neural probabilistic logic programming. \textit{Advances in 
Neural Information Processing Systems}, 31.

\bibitem[Marcus(2020)]{marcus2020next}
Marcus, G. (2020). The next decade in AI: Four steps towards robust artificial 
intelligence. \textit{arXiv preprint arXiv:2002.06177}.

\end{thebibliography}

\end{document}